\documentclass[11pt]{article}
\usepackage[a4paper,margin=22mm]{geometry}
\usepackage{fontspec}
\setmainfont{DejaVu Serif}
\setsansfont{DejaVu Sans}
\usepackage{microtype}
\usepackage{xcolor}
\usepackage{hyperref}
\usepackage{enumitem}
\usepackage{parskip}
\definecolor{Accent}{HTML}{971D32}
\definecolor{Muted}{HTML}{666666}
\hypersetup{colorlinks=true,urlcolor=Accent,linkcolor=Accent}
\setlist{leftmargin=1.6em,itemsep=0.3em,topsep=0.35em}
\setcounter{tocdepth}{2}
\renewcommand{\contentsname}{Contents}
\newcommand{\sectionrule}{\par\vspace{-0.5em}\noindent\textcolor{Accent}{\rule{\linewidth}{0.6pt}}\par}
\begin{document}

\begin{center}
{\sffamily\bfseries\color{Accent} TOEFL WORD OF THE DAY\par}
\vspace{8mm}
{\fontsize{42}{48}\selectfont\bfseries repository\par}
\vspace{2mm}
{\Large\sffamily\color{Accent} /rɪˈpɑːzəˌtɔːri/\par}
\vspace{2mm}
{\small\sffamily Local audio phrase: repository\par}
{\small\sffamily Audio file: audios/repository.mp3\par}
\vspace{2mm}
{\large\sffamily\color{Muted} countable noun\par}
\vspace{5mm}
{\sffamily place or system for storing things \textbullet\ source containing much knowledge\par}
\vspace{6mm}
{\large A repository is a place, system, or source in which valuable information, documents, objects, or knowledge are stored and preserved.\par}
\vspace{6mm}
\begin{minipage}{0.84\textwidth}
\itshape “The university created a digital repository for research papers produced by its faculty.”
\end{minipage}\par
\vspace{10mm}
\textbf{Date:} August 30th 2026\quad\textbullet\quad \textbf{Level:} B1--mid-B2 TOEFL
\vfill
Created and compiled by \textbf{Amir Kooshky}\par
\vspace{2mm}
\href{https://t.me/KooshkyTOEFL}{Telegram Channel}\quad\textbullet\quad
\href{https://www.instagram.com/kooshkytoefl}{Instagram}\quad\textbullet\quad
\href{https://t.me/KooshkyTOEFL\_pv}{Personal Telegram}
\end{center}
\newpage
\tableofcontents
\newpage

\section{Meanings and usage}
\sectionrule
\subsection{Meaning 1: Place or system for storing things}
\textbf{Part of speech:} countable noun

\textbf{IPA:} /rɪˈpɑːzəˌtɔːri/

\textbf{Pronunciation phrase:} repository

\textbf{Local audio file:} audios/repository.mp3

\textbf{Register:} formal; common in academic, scientific, archival, museum, library, and computing contexts

\textbf{Definition:} a place or system where large amounts of information, documents, objects, or other materials are stored and kept

\textbf{In simpler words:} a place where things are stored

\textbf{Typical pattern:} repository for + materials, or repository of + information or objects

\subsubsection{Natural examples}
\begin{enumerate}
  \item The university created a digital repository for research papers produced by its faculty.
  \item The museum serves as a repository for thousands of historical documents and artifacts.
  \item Scientists stored the genetic data in a central online repository.
  \item The national archive acts as a repository for government records of historical importance.
  \item Researchers can download the dataset from a public repository.
  \item The institution maintains a secure repository for confidential documents.
  \item Archaeologists transferred the recovered objects to a repository where they could be properly preserved.
  \item Many universities now maintain institutional repositories containing theses, articles, and other academic work.
\end{enumerate}

\subsubsection{Nuance and implication}
\textbf{Central nuance:} Repository usually suggests organized, purposeful storage rather than simply putting things somewhere temporarily.

\textbf{Usual implication:} The stored material is often valuable, important, or intended to be preserved and accessed later.

\textbf{Compared with temporary storage:} Temporary storage holds something only briefly. A repository is generally intended to preserve and organize material for continued use or safekeeping.

\subsubsection{Grammar notes}
\textbf{How to use it:} Repository is a countable noun. Use a repository for one storage place or system and repositories for more than one.

\textbf{What can follow it:} Use for when emphasizing what the repository is designed to store, and of when describing what it contains.

\textbf{Common preposition:} Both repository for and repository of are common: a repository for research data; a repository of historical records.

\textbf{Noun use:} The singular normally needs a determiner such as a, the, or this.

\subsubsection{Useful patterns}
\textbf{Pattern:} a repository for + material

\textbf{Use:} Use this when emphasizing the purpose of the storage place or system.

\textbf{Example:} The database was designed as a repository for environmental data.

\textbf{Pattern:} a repository of + material

\textbf{Use:} Use this when emphasizing what the repository contains.

\textbf{Example:} The archive is a repository of rare manuscripts and government records.

\textbf{Pattern:} digital repository

\textbf{Use:} Use this for an electronic system used to store and provide access to digital material.

\textbf{Example:} The university's digital repository contains thousands of dissertations.

\textbf{Pattern:} institutional repository

\textbf{Use:} Use this for a system maintained by an institution, especially a university, to preserve its research and publications.

\textbf{Example:} The article is freely available through the university's institutional repository.

\textbf{Pattern:} central repository

\textbf{Use:} Use this for one main location where information or material from different sources is collected.

\textbf{Example:} All laboratory results are uploaded to a central repository.

\subsubsection{Common wrong usage}
\paragraph{Correction 1}
\textbf{Wrong:} The university keeps its papers in a repository place.

\textbf{Correct:} The university keeps its papers in a repository.

\textbf{Why:} Repository already means a place or system for storing things, so place is unnecessary.

\paragraph{Correction 2}
\textbf{Wrong:} The database is a repository to scientific data.

\textbf{Correct:} The database is a repository for scientific data.

\textbf{Why:} Use repository for when describing what a storage system is intended to hold.

\paragraph{Correction 3}
\textbf{Wrong:} The archive contains many repository of documents.

\textbf{Correct:} The archive contains many repositories of documents.

\textbf{Why:} Repository is countable, so use the plural repositories after many.

\subsection{Meaning 2: Source containing much knowledge}
\textbf{Part of speech:} countable noun

\textbf{IPA:} /rɪˈpɑːzəˌtɔːri/

\textbf{Pronunciation phrase:} repository

\textbf{Local audio file:} audios/repository.mp3

\textbf{Register:} formal; especially common in academic, historical, cultural, and literary writing

\textbf{Definition:} a person, book, organization, or other source that contains or preserves a large amount of knowledge or information

\textbf{In simpler words:} a source that holds a lot of knowledge

\textbf{Typical pattern:} a repository of + knowledge, information, memory, tradition, or culture

\subsubsection{Natural examples}
\begin{enumerate}
  \item The archive is an important repository of information about the region's history.
  \item Traditional stories can serve as a repository of a community's cultural knowledge.
  \item For decades, the professor was regarded as a repository of knowledge about ancient languages.
  \item The encyclopedia remains a valuable repository of information for researchers.
  \item Religious texts often function as repositories of a society's values and historical traditions.
  \item The library became a repository of knowledge collected over several centuries.
  \item Older community members may serve as repositories of local history that has never been written down.
\end{enumerate}

\subsubsection{Nuance and implication}
\textbf{Central nuance:} In this figurative use, repository emphasizes that knowledge or cultural information has been accumulated and preserved in a particular source.

\textbf{Usual implication:} The source is valuable because it preserves information that might otherwise be difficult to access or could be lost.

\subsubsection{Grammar notes}
\textbf{How to use it:} Use repository of followed by the kind of knowledge, information, memory, or tradition being preserved.

\textbf{What can follow it:} Common nouns after of include knowledge, information, wisdom, memory, tradition, and culture.

\textbf{Common preposition:} Of is especially common in this meaning: a repository of knowledge.

\textbf{Noun use:} This meaning is countable, so say a repository of knowledge when referring to one source.

\subsubsection{Useful patterns}
\textbf{Pattern:} a repository of knowledge

\textbf{Use:} Use this for a person, institution, text, or collection containing a large amount of accumulated knowledge.

\textbf{Example:} The monastery library became an important repository of knowledge during the medieval period.

\textbf{Pattern:} a repository of information

\textbf{Use:} Use this for a source containing large amounts of useful information.

\textbf{Example:} The database is a valuable repository of information about endangered species.

\textbf{Pattern:} a repository of cultural memory

\textbf{Use:} Use this when something preserves the experiences, traditions, or history of a community.

\textbf{Example:} Oral literature can act as a repository of cultural memory.

\textbf{Pattern:} serve as a repository of + noun

\textbf{Use:} Use this when something performs the function of preserving knowledge or information.

\textbf{Example:} Museums often serve as repositories of cultural knowledge.

\subsubsection{Common wrong usage}
\paragraph{Correction 1}
\textbf{Wrong:} The professor is a repository for knowledge about the subject.

\textbf{Correct:} The professor is a repository of knowledge about the subject.

\textbf{Why:} When describing the knowledge contained in a person or source, repository of is the more natural pattern.

\paragraph{Correction 2}
\textbf{Wrong:} The book is repository of historical information.

\textbf{Correct:} The book is a repository of historical information.

\textbf{Why:} Repository is a singular countable noun, so it normally needs a determiner such as a.

\section{Common collocations}
\sectionrule
\subsection{digital repository}
\textbf{Applies to:} Place or system for storing things

\textbf{Meaning:} an electronic system for storing and providing access to digital materials

\textbf{Example:} The documents are available through a digital repository maintained by the university.

\subsection{institutional repository}
\textbf{Applies to:} Place or system for storing things

\textbf{Meaning:} a digital collection maintained by an institution, especially for its research output

\textbf{Example:} Researchers deposited copies of their articles in the university's institutional repository.

\subsection{data repository}
\textbf{Applies to:} Place or system for storing things

\textbf{Meaning:} a system or location where data is stored and organized

\textbf{Example:} The research team uploaded its results to a public data repository.

\subsection{central repository}
\textbf{Applies to:} Place or system for storing things

\textbf{Meaning:} one main location where material from different sources is stored

\textbf{Example:} The agency created a central repository for records collected by regional offices.

\subsection{public repository}
\textbf{Applies to:} Place or system for storing things

\textbf{Meaning:} a repository that can be accessed by the public

\textbf{Example:} The dataset can be downloaded from a public repository.

\subsection{repository for}
\textbf{Applies to:} Place or system for storing things

\textbf{Meaning:} a place or system intended to store particular materials

\textbf{Pattern:} repository for + material

\textbf{Example:} The archive serves as a repository for official documents.

\subsection{repository of information}
\textbf{Applies to:} Source containing much knowledge

\textbf{Meaning:} a source containing a large amount of information

\textbf{Pattern:} repository of + information

\textbf{Example:} The website has become a major repository of information about the region's wildlife.

\subsection{repository of knowledge}
\textbf{Applies to:} Source containing much knowledge

\textbf{Meaning:} a source in which a large amount of knowledge is preserved

\textbf{Pattern:} repository of + knowledge

\textbf{Example:} The library served as a repository of scientific knowledge.

\subsection{repository of cultural memory}
\textbf{Applies to:} Source containing much knowledge

\textbf{Meaning:} a source that preserves a group's shared history and traditions

\textbf{Example:} Traditional songs can function as a repository of cultural memory.

\section{Synonyms and antonyms}
\sectionrule
\subsection{Close synonyms}
\subsubsection{archive}
\textbf{Part of speech:} countable noun

\textbf{Applies to:} Place or system for storing things

\textbf{Shared meaning:} Both can refer to organized places or systems used to preserve valuable material.

\textbf{Difference:} An archive specifically stores records, documents, or historical materials for long-term preservation. Repository is broader and can contain data, objects, software, documents, or many other kinds of material.

\textbf{Example:} The national archive preserves government documents dating back several centuries.

\subsubsection{database}
\textbf{Part of speech:} countable noun

\textbf{Applies to:} Place or system for storing things

\textbf{Shared meaning:} Both can be electronic systems that store information.

\textbf{Difference:} A database is a structured electronic collection of data designed for searching and processing. A repository is a broader concept and may store files, documents, objects, datasets, or other materials.

\textbf{Example:} The researchers entered the survey responses into a database.

\subsubsection{store}
\textbf{Part of speech:} countable noun

\textbf{Applies to:} Source containing much knowledge

\textbf{Shared meaning:} Both can refer to a large accumulated supply of information or knowledge.

\textbf{Difference:} Store can figuratively mean a supply of knowledge or information, but repository is more formal and emphasizes preservation and accumulation.

\textbf{Example:} The community possesses a rich store of traditional knowledge.

\section{Words learners may confuse}
\sectionrule
\subsection{depository}
\textbf{Part of speech:} countable noun

\textbf{Why learners confuse them:} Both refer to places where things are kept and have similar forms.

\textbf{Key difference:} A repository is any place or system where things are stored and preserved, especially information, documents, or collections. A depository is more specifically a place where something is deposited for safekeeping, especially money, securities, or valuable items.

\textbf{repository:} The university maintains a repository for research data.

\textbf{depository:} The institution acts as a depository for financial securities.

\subsection{archive}
\textbf{Part of speech:} countable noun

\textbf{Why learners confuse them:} Both are frequently used for organized collections intended for long-term preservation.

\textbf{Key difference:} An archive usually preserves historical records or documents. Repository is broader and can store many kinds of material, including digital data and research outputs.

\textbf{repository:} The dataset is stored in an online repository.

\textbf{archive:} The letters are preserved in the university archive.

\subsection{repertory}
\textbf{Part of speech:} countable noun

\textbf{Why learners confuse them:} The words look and sound somewhat similar but have different meanings.

\textbf{Key difference:} Repository means a place or source where things are stored. Repertory refers to a collection of works, skills, or performances available for use, especially in music or theater.

\textbf{repository:} The museum is a repository of historical artifacts.

\textbf{repertory:} The theater company added several new plays to its repertory.

\section{Etymology}
\sectionrule
\textbf{Origin:} Repository comes from Latin repositorium, referring to a place where things are stored or put away.

\textbf{Memory link:} Think of a repository as a place where important material is put away safely so it can be preserved and retrieved later.

\section{Practice exercises}
\sectionrule
\subsection{Correct or incorrect?}
Decide whether each sentence is correct. Correct the inaccurate ones.

\begin{enumerate}
  \item The university maintains a digital repository for research papers.
  \item The database is a repository to scientific data.
  \item The archive serves as a repository of historical documents.
  \item The professor was repository of knowledge about ancient history.
  \item Many universities maintain institutional repositories for academic research.
\end{enumerate}

\subsection{Collocation practice}
Complete each sentence with a natural collocation.

\begin{enumerate}
  \item The university created a digital \_\_\_\_\_ for dissertations and research papers.
  \item The archive serves as a repository \_\_\_\_\_ thousands of historical records.
  \item The library became an important repository \_\_\_\_\_ scientific knowledge.
  \item All experimental data are uploaded to a \_\_\_\_\_ repository maintained by the research institute.
\end{enumerate}

\subsection{Complete answer key}
\subsubsection{Correct or incorrect?}
\begin{enumerate}
  \item \textbf{Correct} \emph{A digital repository is an electronic system for storing and preserving digital material.}
  \item \textbf{Incorrect}. The database is a repository for scientific data. \emph{Use repository for when describing what a storage system is intended to contain.}
  \item \textbf{Correct} \emph{Repository of naturally describes the material contained in a repository.}
  \item \textbf{Incorrect}. The professor was a repository of knowledge about ancient history. \emph{Repository is a singular countable noun and normally requires a determiner such as a.}
  \item \textbf{Correct} \emph{Institutional repository is a standard term for a system that preserves an institution's research output.}
\end{enumerate}

\subsubsection{Collocation practice}
\begin{enumerate}
  \item \textbf{repository.} The university created a digital repository for dissertations and research papers. \emph{A digital repository stores electronic documents and other digital materials.}
  \item \textbf{for.} The archive serves as a repository for thousands of historical records. \emph{Repository for emphasizes what the repository is intended to store.}
  \item \textbf{of.} The library became an important repository of scientific knowledge. \emph{Repository of emphasizes what knowledge or information the source contains.}
  \item \textbf{central.} All experimental data are uploaded to a central repository maintained by the research institute. \emph{A central repository is one main location where information from different sources is collected.}
\end{enumerate}

\vfill
\begin{center}
\small\color{Muted} Created and compiled by \textbf{Amir Kooshky}\\
\href{https://t.me/KooshkyTOEFL}{Telegram Channel}\quad\textbullet\quad
\href{https://www.instagram.com/kooshkytoefl}{Instagram}\quad\textbullet\quad
\href{https://t.me/KooshkyTOEFL\_pv}{Personal Telegram}
\end{center}
\end{document}
